\section{A spurious feature as a field shift}
\label{sec:spurious}

We now perturb the encoding, and only the encoding.

A fraction $f_d$ of agents, drawn independently of their own label, extend their
representation of an incoming message with a feature that carries no information
about the issue and is determined entirely by the \emph{sender's} label. This is
the shortcut-learning setup moved from one classifier to a population: a feature
that is predictive of nothing the task cares about, but perfectly predictive of a
group membership. Because it enters the receiver's encoding of the message, its
effect on Eq.~\ref{eq:fields} is an additive shift of the opinion field,
\begin{equation}
  \hw \;\longrightarrow\; \hw + D\bigl[\kappa_r, \kappa_e\bigr],
  \label{eq:shift}
\end{equation}
with $D$ a $2 \times 2$ matrix indexed by the receiver's and the emitter's labels.
The main text uses the fully symmetric case, in which both groups are affected
alike:
\begin{equation}
  D \;=\; d \cdot \begin{pmatrix} +1 & -1 \\ -1 & +1 \end{pmatrix},
  \qquad d > 0 .
  \label{eq:dmatrix}
\end{equation}
The five asymmetric variants --- one group only, one direction only, and so on ---
are implemented and reported in Appendix~\ref{app:variants}.

\paragraph{Sign convention.} With Eq.~\ref{eq:dmatrix} and $d > 0$, a same-label
sender's message arrives with $\hw$ shifted \emph{up}, toward agreement, and an
opposite-label sender's shifted \emph{down}, toward disagreement. Since $F_\mu$
carries the prefactor $(1 - 2\Phi(\hw))$, agreement builds trust and disagreement
erodes it, so $d > 0$ is \textbf{in-group tolerance and out-group intolerance} and
$d < 0$ is its reverse. We flag this explicitly because the sign is easy to get
backwards, and because it fixes the orientation of the $x$ axis in every figure
that follows.

\paragraph{What has and has not been assumed.} No training data is biased. No
agent holds a preference over groups. The label is uncorrelated with every issue
on the agenda, and $f_d$ is independent of label, so neither group is
over-represented among the perturbed agents. A discriminating agent runs
\emph{the same update} as every other agent; it has merely encoded the message
badly, in a way that happens to depend on who sent it. Everything in
\S\ref{sec:phase} is downstream of Eq.~\ref{eq:shift} alone.

\paragraph{Why this should be harmless, and is not.} An uninformative feature
appended to an input is, for a single learner, just noise: it costs sample
efficiency and nothing else. That intuition is correct here too, and we verify it
--- keying the same shift to a per-source \emph{random} bit rather than the label
degrades learning and produces no group structure (\S\ref{sec:ablation}, A0). The
difference is that the trust sector is indexed by source. A shift that is constant
across a source's messages accumulates in $\mu_{e|r}$ rather than averaging out,
and once $\mu$ is class-correlated the anti-learning prefactor of
\S\ref{sec:rule} makes the opinion sector class-correlated too. The feedback runs
in both directions, which is what produces a threshold rather than a gradual
degradation.
